\subsection{W 在 T 上的线性表示的分解}

令 I 为 W 的 Coxeter 图的顶点集（第 4 号），令 J 为 I 的一个子集，使 J 中没有顶点与 I-J 中的任何顶点相连。设 C 是一个室，s 是从 I 到关于 C 的墙的反射集合的规范双射，并令 \(W_{J,C}\) 为由像 \(s(\mathrm{J})\) 生成的子群。由《代数》第 IV 章 §1 第 9 号命题 8 可知，W 是两个子群 \(W_{J,C}\) 和 \(W_{I-J,C}\) 的直积。设 \(C'\) 是另一个室，\(s'\) 是从 I 到 W 的相应单射。我们已经看到（第 4 号），若 \(w \in W\) 将 C 变换为 \(C'\)，则有 \(s'(\mathrm{i}) = ws(\mathrm{i})w^{-1}\)，其中 \(i \in I\)。由于 \(W_{J,C}\) 在 W 中正规，可知 \(s'(\mathrm{i}) \in W_{J,C}\) 对所有 \(i \in J\) 成立。我们推得子群 \(W_{J,C}\) 与 C 无关。从现在起将其简记为 \(W_{J}\)。

\(W_{J,C}\) 的定义显然可推广到 I 的任意子集 J。但是，如果 J 中有一个顶点与 I-J 中的一个顶点相连，则 \(W_{J,C}\) 不正规并且取决于 C 的选择。

令 \(T_{J}^{0}\) 为 T 中在 \(U(W_{J})\) 的每个元素下不变的向量所成的子空间，令 \(T_{J}\) 为与 \(T_{J}^{0}\) 正交的子空间。由于 \(W_{J}\) 是 W 的正规子群，显然 \(T_{J}^{0}\) 在 \(U(W)\) 下不变，因而 \(T_{J}\) 也是如此。

命题 5. 设 \(J_{1},\ldots,J_{s}\) 是 W 的 Coxeter 图的连通分支的顶点集。对于 \(1\leqslant p\leqslant s\)，置

\[
\mathrm{W} _ {p} = \mathrm{W} _ {\mathrm{J} _ {p}}, \quad \mathrm{T} _ {p} = \mathrm{T} _ {\mathrm{J} _ {p}}; \quad \mathrm{T} _ {p} ^ {\prime} = \mathrm{T} _ {\mathrm{J} _ {p}} ^ {0}, \quad \text {and} \quad \mathrm{T} _ {0} = \bigcap_ {1 \leqslant p \leqslant s} \mathrm{T} _ {p} ^ {\prime}.
\]

\begin{enumerate}
\def\labelenumi{(\roman{enumi})}
\item
群 W 是群 \(\mathrm{W}_p\) 的直积（\(1 \leqslant p \leqslant s\)）。
\item
空间 \(\mathbf{T}\) 是子空间 \(\mathrm{T}_0, \mathrm{T}_1, \ldots, \mathrm{T}_s\) 的正交直和，它们都在 \(\mathrm{U}(\mathrm{W})\) 下稳定。
\item
对于所有 \(q\) 满足 \(1 \leqslant q \leqslant s\)，子空间 \(\mathrm{T}_q'\) 是 \(\mathrm{T}\) 中由在 \(\mathrm{U}(\mathrm{W}_q)\) 的每个元素下不变的向量所成的；它是 \(\mathrm{T}_p\)（\(1 \leqslant p \leqslant s\) 且 \(p \neq q\)）的直和。
\item
设 \(\mathbf{C}\) 是一个室。子空间 \(\mathrm{T}_p\)（\(1 \leqslant p \leqslant s\)）由向量 \(e_i(\mathbf{C})\)（其中 \(i \in \mathrm{J}_p\)）生成（按第 4 号的记号）。
\item
W 在子空间 \(T_{p}\)（\(1 \leqslant p \leqslant s\)）上的表示是绝对不可约、非平凡且两两不等价的。
\end{enumerate}

断言 (i) 由《代数》第 IV 章 §1 第 9 号得出。另一方面，我们已经看到子空间 \(\mathrm{T}_p\) 在 \(\mathrm{U}(\mathrm{W})\) 下不变，\(\mathrm{T}_0\) 也如此。设 \(\mathrm{C}\) 是一个室；由于 \(\mathrm{W}_p\) 由反射 \(s_i(\mathrm{C})\)（其中 \(i \in \mathrm{J}_p\)）生成，显然 \(\mathrm{T}_p'\) 是与 \(e_i(\mathrm{C})\)（其中 \(i \in \mathrm{J}_p\)）正交的子空间，因而得到 (iv)。此外，若 \(i \in \mathrm{J}_p\)、\(j \in \mathrm{J}_q\) 且 \(p \neq q\)，则 \(m_{ij} = 2\)，因为 \(\{i,j\}\) 不是 \(\mathrm{W}\) 的 Coxeter 图的一条边，所以 \((e_i|e_j) = 0\)。断言 (ii) 现已显然。断言 (iii) 则随之得出，因为 \(\mathrm{T}_q'\) 是 \(\mathrm{T}_q\) 的正交补。

最后，设 V 是在 \(T_{p}\) 下不变的 \(\mathrm{U}(\mathrm{W}_{p})\) 的子空间。对于所有 \(i \in J_{p}\)，要么 \(e_{i} \in V\)，要么 \(e_{i}\) 与 V 正交（\(\S2\)2，第 2 号，命题 3）。令 A（分别为 B）为满足 \(i \in J_{p}\)（分别为 \(e_{i} \in V\) 与 V 正交）的 \(e_{i}\) 所成的集合。显然，对 \((e_{i}|e_{j}) = 0\) 且 \(i \in A\) 有 \(j \in B\)，而由于 \(J_{p}\) 连通，可知或者 A = \(\varnothing\) 且 V = \(\{0\}\)，或者 A = \(J_{p}\) 且 V = \(T_{p}\)。因此，\(W_{p}\) 在 \(T_{p}\) 上的表示是不可约的，从而是绝对不可约的（\(\S2\)2，第 1 号，命题 1）。根据 \(T_{p}\) 的定义，它是非平凡的。最后，(v) 中的最后一个断言立即由 (iii) 得出。

如果 T 中在 U(W) 下不变的子空间 \(T_{0}\) 缩减为 \(\{0\}\)，则称 W 是本质的；如果 W 在 T 上的表示 U 是不可约的，则称 W 是不可约的。

推论. 假设 \(W \neq \{1\}\)。则 W 不可约，当且仅当它是本质的且其 Coxeter 图连通。

注. 在命题 5 的假设下，\(T_{p}\)（\(0 \leqslant p \leqslant s\)）是 W 在 T 上的线性表示 U 的等型分量（《代数》第 VIII 章 §3 第 4 号）。由此（同上，命题 11），T 中每个在 W 下稳定的向量子空间 V 都是 \(V \cap T_{p}\)（\(0 \leqslant p \leqslant s\)）之和；此外（同上，命题 10），每个与 \(U(w)\) 的算子 \(w \in W\) 对易的自同态都使每个 \(T_{p}\)（\(0 \leqslant p \leqslant s\)）保持稳定，并且对 \(1 \leqslant p \leqslant s\) 在其上诱导一个仿射伸缩。特别地，T 上在 W 下不变的双线性型 \(\Phi\) 是由下式给出的双线性型

\[
\varPhi \left(\sum_ {k} t _ {k}, \sum_ {k} t _ {k} ^ {\prime}\right) = \varPhi_ {0} (t _ {0}, t _ {0} ^ {\prime}) + \sum_ {1 \leqslant p \leqslant s} a _ {p} (t _ {p} | t _ {p} ^ {\prime}),
\]

其中 \(\Phi_{0}\) 是 \(T_{0}\) 上的双线性型，而 \(a_{p}\)（对 \(1 \leqslant p \leqslant s\)）是一个实数：事实上，这样的型 \(\Phi\) 可以唯一地写成 \((t, t') \mapsto (\mathrm{A}(t)|t')\) 的形式，其中 A 是一个与 \(\mathrm{U}(w)\) 的 \(w \in W\) 对易的 T 的自同态。
% semantic-fragments: legacy-input-seam
\relax{}
